\documentclass[11pt,a4paper]{article}
\usepackage[utf8]{inputenc}
\usepackage[T1]{fontenc}
\usepackage{lmodern}
\usepackage{geometry}
\usepackage{graphicx}
\usepackage{booktabs}
\usepackage{array}
\usepackage{longtable}
\usepackage{hyperref}
\usepackage{xcolor}
\usepackage{float}
\usepackage{enumitem}
\usepackage{fancyhdr}
\usepackage{titlesec}
\usepackage{multicol}

\geometry{margin=2.5cm}

% Colors
\definecolor{primaryblue}{RGB}{0,82,147}
\definecolor{secondaryblue}{RGB}{0,114,189}
\definecolor{accentgreen}{RGB}{0,150,100}

% Hyperlink setup
\hypersetup{
    colorlinks=true,
    linkcolor=primaryblue,
    citecolor=secondaryblue,
    urlcolor=accentgreen
}

% Section formatting
\titleformat{\section}{\Large\bfseries\color{primaryblue}}{\thesection}{1em}{}
\titleformat{\subsection}{\large\bfseries\color{secondaryblue}}{\thesubsection}{1em}{}
\titleformat{\subsubsection}{\normalsize\bfseries\color{accentgreen}}{\thesubsubsection}{1em}{}

\pagestyle{fancy}
\fancyhf{}
\fancyhead[L]{\small Bibliographic Report on Security in IaC and Smart Contracts}
\fancyhead[R]{\small\thepage}
\renewcommand{\headrulewidth}{0.4pt}

\title{\textbf{Bibliographic Report on Security Analysis\\in Infrastructure-as-Code and Smart Contracts}\\[0.5em]
\large A Comprehensive Survey of Models, Methods, and Proposals}
\author{Research Analysis}
\date{\today}

\begin{document}

\maketitle

\begin{abstract}
This bibliographic report presents a comprehensive analysis of ten research papers focusing on security aspects in Infrastructure-as-Code (IaC) and Smart Contracts. The papers span from 2022 to 2025, representing the state-of-the-art in vulnerability detection, security smell taxonomy, and the application of Large Language Models (LLMs) for security analysis. We summarize the key models, methodologies, and findings from each paper, followed by a synthesis and proposal for future research directions. The analysis reveals a clear trend toward LLM-based approaches for security analysis, while highlighting the continued importance of traditional static analysis tools. We identify gaps in current research and propose an integrated framework that combines the strengths of both approaches.
\end{abstract}

\tableofcontents
\newpage

%==============================================================================
\section{Introduction}
%==============================================================================

Infrastructure-as-Code (IaC) has become a cornerstone of modern DevOps practice, enabling the automated provisioning and management of IT infrastructure through code and automation. Similarly, smart contracts on blockchain platforms like Ethereum have revolutionized decentralized application development. However, both domains face significant security challenges that can lead to severe financial and operational consequences.

This bibliographic report analyzes ten research papers that address security concerns in IaC and smart contracts from various perspectives. The papers can be broadly categorized into three main themes: (1) vulnerability detection and taxonomy development, (2) application of Large Language Models (LLMs) for security analysis, and (3) comparative evaluation of detection tools and methods.

The evolution of research in this field shows a clear trajectory from rule-based static analysis tools toward more sophisticated approaches leveraging machine learning and LLMs. This report synthesizes the findings from these papers and provides insights into the current state-of-the-art, remaining challenges, and potential future directions.

%==============================================================================
\section{Summary of Analyzed Papers}
%==============================================================================

\subsection{Paper 1: Vulnerabilities in Infrastructure as Code (EMSE 2025)}

\textbf{Authors:} Aicha War, Alioune Diallo, Andrew Habib, Jacques Klein, Tegawendé F. Bissyandé

\textbf{Publication:} Empirical Software Engineering, 2025

\subsubsection{Overview}
This comprehensive empirical study investigates security vulnerabilities across Infrastructure-as-Code components, including tools, scripts, and add-ons. The research employs Static Application Security Testing (SAST) tools, specifically Snyk and Horusec, to analyze over 1,600 repositories encompassing seven widely adopted IaC tools including Ansible, Terraform, Chef, Puppet, Saltstack, Pulumi, and Vagrant.

\subsubsection{Key Findings}
The study reveals that IaC components are plagued by exploitable vulnerabilities spanning all ten categories of the OWASP Top 10 2021. A particularly significant finding is that manifest files emerge as the most vulnerable artifacts within the IaC ecosystem. The research also uncovers that IaC vulnerabilities frequently stem from non-specialized actors attempting modifications within components outside their expertise, highlighting the importance of cross-functional collaboration in DevSecOps practices.

\subsubsection{Methodology}
The methodology relies on widely adopted static security testing tools to analyze over 1,600 repositories. The researchers employed Jenkins as a CI tool for automation, combined with analysis reports from Snyk and Horusec. The categorization process followed the OWASP Top 10 2021 framework, providing a standardized approach to vulnerability classification.

\subsubsection{Models and Tools}
\begin{itemize}
    \item \textbf{SAST Tools:} Snyk, Horusec
    \item \textbf{Classification Framework:} OWASP Top 10 2021
    \item \textbf{Automation:} Jenkins CI/CD pipeline
    \item \textbf{Data Sources:} 1,600+ GitHub repositories
\end{itemize}

%------------------------------------------------------------------------------
\subsection{Paper 2: GLITCH - Automated Polyglot Security Smell Detection (ASE 2022)}

\textbf{Authors:} Nuno Saavedra, João F. Ferreira

\textbf{Publication:} IEEE/ACM International Conference on Automated Software Engineering (ASE), 2022

\subsubsection{Overview}
GLITCH represents a significant advancement in security smell detection for IaC scripts by introducing a technology-agnostic framework that enables automated polyglot detection. The key innovation lies in transforming IaC scripts into an intermediate representation, upon which different security smell detectors can be defined uniformly across different IaC technologies.

\subsubsection{Key Findings}
The large-scale empirical study analyzed security smells across three datasets containing 196,755 IaC scripts and 12,281,251 lines of code. All categories of security smells were identified across all datasets, with hard-coded secrets and suspicious comments being the most dominant security smells across Ansible, Chef, and Puppet scripts.

\subsubsection{Methodology}
GLITCH employs an intermediate representation that can model scripts written in different IaC technologies. The framework defines nine rules that operate on this intermediate representation to detect security smells. The evaluation compared GLITCH with state-of-the-art security smell detectors (SLIC and SLAC), demonstrating higher precision and recall.

\subsubsection{Models and Tools}
\begin{itemize}
    \item \textbf{Framework:} GLITCH (technology-agnostic intermediate representation)
    \item \textbf{Supported Languages:} Ansible, Chef, Puppet
    \item \textbf{Detection Rules:} 9 security smell categories
    \item \textbf{Comparison Baselines:} SLIC, SLAC
\end{itemize}

%------------------------------------------------------------------------------
\subsection{Paper 3: DAppSCAN - Smart Contract Weakness Datasets (TSE 2024)}

\textbf{Authors:} Zibin Zheng, Jianzhong Su, Jiachi Chen, David Lo, Zhijie Zhong, Mingxi Ye

\textbf{Publication:} IEEE Transactions on Software Engineering, June 2024

\subsubsection{Overview}
DAppSCAN addresses the critical need for large-scale, unbiased, real-world datasets for smart contract weakness detection. The research involved 22 participants spending 44 person-months analyzing 1,199 open-source audit reports from 29 security teams, resulting in the identification of 9,154 weaknesses.

\subsubsection{Key Findings}
Two distinct datasets were developed: DAPPSCAN-SOURCE comprising 39,904 Solidity files with 1,618 SWC weaknesses from 682 real-world DApps, and DAPPSCAN-BYTECODE containing 6,665 compiled smart contracts with 888 SWC weaknesses. The evaluation of state-of-the-art smart contract weakness detection tools revealed sub-par performance in terms of both effectiveness and success detection rate.

\subsubsection{Methodology}
The researchers manually analyzed open-source audit reports to identify and categorize weaknesses. A tool was developed to automatically identify dependency relationships within DApp projects and complete missing public libraries, enabling the creation of compilable contract datasets.

\subsubsection{Models and Tools}
\begin{itemize}
    \item \textbf{Datasets:} DAPPSCAN-SOURCE, DAPPSCAN-BYTECODE
    \item \textbf{Classification:} Smart Contract Weakness Classification (SWC) Registry
    \item \textbf{Evaluated Tools:} Mythril, Securify, Slither, Smartian, Sailfish, eTainter, Echidna
    \item \textbf{Scale:} 1,199 audit reports, 9,154 weaknesses
\end{itemize}

%------------------------------------------------------------------------------
\subsection{Paper 4: Defect Taxonomy for IaC - A Replication Study (arXiv 2025)}

\textbf{Authors:} Wendell Oliveira, Filipe Paiva, Thiago Emmanuel Pereira, João Brunet

\textbf{Publication:} arXiv:2505.01568v3, 2025

\subsubsection{Overview}
This replication study validates the "Gang of Eight" defect taxonomy proposed by Rahman et al. across Programming Language-based IaC (PL-IaC) tools including Terraform, Pulumi, and AWS CDK. The research extends the original study by examining both open-source and proprietary projects from VTEX and Nubank.

\subsubsection{Key Findings}
The research confirmed the same eight defect categories identified in the original study: Conditional, Configuration Data, Dependency, Documentation, Idempotency, Security, Service, and Syntax. Configuration Data defects maintain high frequency in both open-source and proprietary codebases, appearing in over 30\% of Nubank's PL-IaC programs.

\subsubsection{Methodology}
The study performed qualitative analysis on 3,364 defect-related commits from 285 open-source PL-IaC repositories (PIPr dataset). The ACID tool was enhanced to automatically classify and analyze defect distributions across an expanded dataset including proprietary projects.

\subsubsection{Models and Tools}
\begin{itemize}
    \item \textbf{Taxonomy:} "Gang of Eight" defect categories
    \item \textbf{Detection Tool:} ACID (enhanced version)
    \item \textbf{PL-IaC Tools:} Terraform, Pulumi, AWS CDK
    \item \textbf{Datasets:} PIPr, VTEX, Nubank
\end{itemize}

%------------------------------------------------------------------------------
\subsection{Paper 5: Forge - LLM-Driven Vulnerability Dataset Construction (ICSE 2026)}

\textbf{Authors:} Jiachi Chen, Yiming Shen, Jiashuo Zhang, et al.

\textbf{Publication:} IEEE/ACM International Conference on Software Engineering (ICSE), 2026

\subsubsection{Overview}
Forge introduces the first automated framework for constructing comprehensive smart contract vulnerability datasets using Large Language Models. The framework leverages an LLM-driven pipeline to extract high-quality vulnerabilities from real-world audit reports and classify them according to the Common Weakness Enumeration (CWE).

\subsubsection{Key Findings}
Running Forge on 6,454 real-world audit reports generated a dataset comprising 81,390 Solidity files and 27,497 vulnerability findings across 296 CWE categories. The framework achieved a Krippendorff's alpha coefficient of 0.87, indicating substantial agreement between Forge and human experts in CWE classification.

\subsubsection{Methodology}
Forge employs a divide-and-conquer strategy through a map-reduce paradigm to extract structured and self-contained vulnerability information from unstructured audit reports. A tree-of-thoughts technique is used to classify vulnerability information into the hierarchical CWE classification system.

\subsubsection{Models and Tools}
\begin{itemize}
    \item \textbf{LLM:} GPT-based models for extraction and classification
    \item \textbf{Classification:} Common Weakness Enumeration (CWE)
    \item \textbf{Key Components:} Semantic Chunker, MapReduce Extractor, Hierarchical Classifier, Code Fetcher
    \item \textbf{Dataset Scale:} 81,390 Solidity files, 27,497 vulnerabilities
\end{itemize}

%------------------------------------------------------------------------------
\subsection{Paper 6: LLMs vs Static Code Analysis Tools - A Benchmark (IEEE Access 2025)}

\textbf{Authors:} Damian Gnieciak, Tomasz Szandala

\textbf{Publication:} IEEE Access, 2025

\subsubsection{Overview}
This study provides a systematic comparison between three industry-standard rule-based static code-analysis tools (SonarQube, CodeQL, and Snyk Code) and three state-of-the-art large language models (GPT-4.1, Mistral Large, and DeepSeek V3) for vulnerability detection.

\subsubsection{Key Findings}
The language-based scanners achieved higher mean F-1 scores (0.797, 0.753, and 0.750 for LLMs) than their static counterparts (0.260, 0.386, and 0.546). LLMs' advantage originates from superior recall, confirming their ability to reason across broader code contexts. However, LLMs exhibit higher false-positive ratios and struggle with precise issue localization due to tokenization artifacts.

\subsubsection{Methodology}
The evaluation used a curated suite of ten real-world C\# projects embedding 63 vulnerabilities across categories such as SQL injection, hard-coded secrets, and outdated dependencies. Classical detection accuracy metrics (precision, recall, F-score), analysis latency, and developer effort were measured.

\subsubsection{Models and Tools}
\begin{itemize}
    \item \textbf{Static Analysis Tools:} SonarQube, CodeQL, Snyk Code
    \item \textbf{LLMs:} GPT-4.1, Mistral Large, DeepSeek V3
    \item \textbf{Evaluation Metrics:} Precision, Recall, F1-Score
    \item \textbf{Test Dataset:} 10 C\# projects with 63 vulnerabilities
\end{itemize}

%------------------------------------------------------------------------------
\subsection{Paper 7: Security Smells Taxonomy Update Beyond Seven Sins (2025)}

\textbf{Authors:} Aicha War, Serge L.B. Nikiema, Jordan Samhi, Jacques Klein, Tegawendé F. Bissyandé

\textbf{Publication:} arXiv:2509.18761v1, 2025

\subsubsection{Overview}
This paper revisits the security smell taxonomy by extending the study to seven popular IaC tools (Ansible, Terraform, Chef, Puppet, Pulumi, Saltstack, and Vagrant) and incorporating automation through LLMs. The study yields a comprehensive taxonomy of 62 security smell categories, significantly expanding beyond the previously known seven.

\subsubsection{Key Findings}
The top 10 most represented security smell categories include Insecure Configuration Management, Insecure Dependency Management, Insecure Input Handling, Outdated Dependencies, Path Traversal, Command Injection, Code Injection, Outdated Software Version, Inadequate Naming Convention, and Sensitive Information Exposure. The study demonstrates that security issues persist for extended periods, likely due to inadequate detection and mitigation tools.

\subsubsection{Methodology}
The research employed LLMs (GPT-3.5-Turbo and GPT-4o) for initial pattern processing with systematic human validation and reconciliation with established security standards (CWE). New security checking rules were implemented within linters for seven popular IaC tools, often achieving 1.00 precision score.

\subsubsection{Models and Tools}
\begin{itemize}
    \item \textbf{LLMs:} GPT-3.5-Turbo, GPT-4o
    \item \textbf{Taxonomy Size:} 62 security smell categories
    \item \textbf{IaC Tools:} Ansible, Terraform, Chef, Puppet, Pulumi, Saltstack, Vagrant
    \item \textbf{Linter Enhancement:} Ansible-lint, Rubocop, Terrascan, Salt-Lint, ESLint, Bandit, Yaml-Lint
\end{itemize}

%------------------------------------------------------------------------------
\subsection{Paper 8: Semantics-Aware Detection of Security Smells in IaC (2025)}

\textbf{Authors:} Aicha War, Adnan A. Rawass, Abdoul K. Kabore, Jordan Samhi, Jacques Klein, Tegawendé F. Bissyandé

\textbf{Publication:} arXiv:2509.18790v1, 2025

\subsubsection{Overview}
This work introduces a novel approach that enhances static analysis with semantic understanding by jointly leveraging natural language and code representations. The method builds on two complementary ML models: CodeBERT to capture semantics across code and text, and LongFormer to represent long IaC scripts without losing contextual information.

\subsubsection{Key Findings}
Semantic enrichment substantially improves detection, raising precision and recall from 0.46 and 0.79 to 0.92 and 0.88 on Ansible, and from 0.55 and 0.97 to 0.87 and 0.75 on Puppet, respectively. The ablation study confirmed that removing semantic context significantly degrades model performance.

\subsubsection{Methodology}
The approach was evaluated on misconfiguration datasets from two widely used IaC tools (Ansible and Puppet). Two ablation studies were conducted: removing code text from natural language input, and truncating scripts to reduce context. The method was compared against four LLMs and prior work.

\subsubsection{Models and Tools}
\begin{itemize}
    \item \textbf{ML Models:} CodeBERT, LongFormer
    \item \textbf{Comparison LLMs:} GPT-4-Turbo, GPT-3.5-Turbo, Starcoder-2, LLaMA-3.2
    \item \textbf{Target Tools:} Ansible, Puppet
    \item \textbf{Key Innovation:} Joint natural language and code processing
\end{itemize}

%------------------------------------------------------------------------------
\subsection{Paper 9: GenSIaC - Security-Aware IaC Generation with LLMs (2025)}

\textbf{Authors:} Yikun Li, Matteo Grella, Daniel Nahmias, et al.

\textbf{Publication:} arXiv:2511.12385v1, 2025

\subsubsection{Overview}
GenSIaC investigates the potential of Large Language Models in generating security-aware IaC code. The paper proposes an instruction fine-tuning dataset designed to improve LLMs' ability to recognize potential security weaknesses during code generation and inspection.

\subsubsection{Key Findings}
Code-specific LLMs showed poor F1-scores (0.000-0.303) for security weakness recognition, while open-source general-purpose models (Vicuna and Llama2) performed better (0.354-0.495). After fine-tuning CodeLlama with GenSIaC, F1-scores improved from 0.276-0.303 to 0.771-0.858, surpassing larger LLMs without fine-tuning.

\subsubsection{Methodology}
The study first evaluated non-fine-tuned LLMs in recognizing security misconfigurations during generation and inspection of IaC code. GenSIaC, a fine-tuning dataset, was then introduced to enhance LLMs' security weakness awareness in IaC scripts. Evaluation covered nine IaC security weakness categories.

\subsubsection{Models and Tools}
\begin{itemize}
    \item \textbf{Base Models:} CodeGen2.5, StarCoder, WizardCoder, CodeLlama, Vicuna, Llama2, GPT-3.5, GPT-4
    \item \textbf{Fine-tuning Dataset:} GenSIaC
    \item \textbf{Security Categories:} 9 IaC security weakness types
    \item \textbf{Key Metric Improvement:} F1-score from 0.276 to 0.771+
\end{itemize}

%------------------------------------------------------------------------------
\subsection{Paper 10: Agent READMEs - Context Files for Agentic Coding (2025)}

\textbf{Authors:} Worawalan Chatlatanagulchai, Hao Li, Yutaro Kashiwa, et al.

\textbf{Publication:} arXiv:2511.12884v1, 2025

\subsubsection{Overview}
This paper conducts the first large-scale empirical study of 2,303 agent context files from 1,925 repositories to characterize their structure, maintenance, and content. Agent context files serve as persistent, project-level instructions for AI coding agents.

\subsubsection{Key Findings}
Developers prioritize functional context such as build and run commands (62.3\%), implementation details (69.9\%), and architecture (67.7\%). A significant gap was identified: non-functional requirements like security (14.5\%) and performance (14.5\%) are rarely specified. Agent context files are actively maintained in short bursts, evolving through incremental additions rather than deletions.

\subsubsection{Methodology}
The study analyzed context files for three agentic coding tools: ClaudeCode (CLAUDE.md), OpenAI Codex (AGENTS.md), and GitHub Copilot (copilot-instructions.md). Content analysis identified 16 instruction types, and automatic classification achieved 0.79 F1-score for concrete functional topics.

\subsubsection{Models and Tools}
\begin{itemize}
    \item \textbf{Agentic Tools:} ClaudeCode, OpenAI Codex, GitHub Copilot
    \item \textbf{Context Files:} CLAUDE.md, AGENTS.md, copilot-instructions.md
    \item \textbf{Dataset Scale:} 2,303 files from 1,925 repositories
    \item \textbf{Instruction Categories:} 16 types identified
\end{itemize}

%==============================================================================
\section{Comparative Analysis of Models and Methods}
%==============================================================================

\subsection{Evolution of Detection Approaches}

The analyzed papers reveal a clear evolution in security detection approaches for IaC and smart contracts:

\begin{enumerate}
    \item \textbf{Rule-Based Static Analysis (2022-2024):} Early approaches relied on predefined patterns and rules, exemplified by tools like SLIC, SLAC, and traditional SAST tools. These methods achieve high precision but often miss novel or context-dependent vulnerabilities.
    
    \item \textbf{Intermediate Representations (2022):} GLITCH introduced the concept of technology-agnostic intermediate representations, enabling consistent detection across multiple IaC languages while maintaining reasonable accuracy.
    
    \item \textbf{Machine Learning Approaches (2024-2025):} Papers 8 and 9 demonstrate the effectiveness of ML models like CodeBERT and LongFormer for semantic-aware security detection, achieving significant improvements over traditional approaches.
    
    \item \textbf{LLM-Based Analysis (2025-2026):} The most recent papers (5, 6, 7, 9) showcase the application of LLMs for both security detection and secure code generation, with fine-tuning approaches achieving state-of-the-art performance.
\end{enumerate}

\subsection{Performance Comparison}

Table \ref{tab:performance} summarizes the key performance metrics reported across the analyzed papers:

\begin{table}[H]
\centering
\caption{Performance Comparison of Detection Approaches}
\label{tab:performance}
\begin{tabular}{llccc}
\toprule
\textbf{Paper} & \textbf{Method} & \textbf{Precision} & \textbf{Recall} & \textbf{F1-Score} \\
\midrule
Paper 2 & GLITCH & 0.77 & 0.87 & 0.82 \\
Paper 6 & SonarQube & 0.57 & 0.19 & 0.26 \\
Paper 6 & CodeQL & 0.63 & 0.28 & 0.39 \\
Paper 6 & Snyk Code & 0.69 & 0.52 & 0.55 \\
Paper 6 & GPT-4.1 & 0.76 & 0.88 & 0.80 \\
Paper 6 & Mistral Large & 0.78 & 0.78 & 0.75 \\
Paper 6 & DeepSeek V3 & 0.72 & 0.85 & 0.75 \\
Paper 8 & CodeBERT (Ansible) & 0.92 & 0.88 & 0.90 \\
Paper 8 & LongFormer (Puppet) & 0.87 & 0.75 & 0.79 \\
Paper 9 & CodeLlama (fine-tuned) & 0.77+ & -- & 0.77-0.86 \\
\bottomrule
\end{tabular}
\end{table}

\subsection{Taxonomy Evolution}

The security smell taxonomies have evolved significantly:

\begin{itemize}
    \item \textbf{Original Seven Sins (2019):} Rahman et al. identified 7 security smell categories for Puppet scripts
    \item \textbf{Gang of Eight (2020):} Expanded to 8 defect categories across multiple dimensions
    \item \textbf{GLITCH Nine (2022):} Defined 9 security smell types for polyglot detection
    \item \textbf{Extended 62 Categories (2025):} Paper 7 significantly expanded to 62 fine-grained security smell categories
\end{itemize}

%==============================================================================
\section{Proposed Framework}
%==============================================================================

Based on the comprehensive analysis of the ten papers, we propose an integrated framework for security analysis in IaC and smart contracts that combines the strengths of multiple approaches.

\subsection{Multi-Tier Security Analysis Pipeline}

We propose a three-tier architecture for comprehensive security analysis:

\subsubsection{Tier 1: Fast Pattern Matching}
The first tier employs rule-based static analysis for quick identification of known vulnerability patterns. This tier leverages intermediate representations (as demonstrated by GLITCH) for technology-agnostic detection across multiple IaC languages. Key components include:
\begin{itemize}
    \item Pre-defined security rules based on CWE and OWASP classifications
    \item Intermediate representation transformation for polyglot support
    \item Low-latency detection suitable for CI/CD integration
\end{itemize}

\subsubsection{Tier 2: Semantic-Aware ML Detection}
The second tier employs fine-tuned language models for deeper semantic analysis:
\begin{itemize}
    \item CodeBERT-based analysis for capturing code-text relationships
    \item LongFormer for handling long-context scripts
    \item Integration of natural language context (comments, documentation) for improved accuracy
\end{itemize}

\subsubsection{Tier 3: LLM-Assisted Deep Analysis}
The third tier leverages LLMs for comprehensive security review:
\begin{itemize}
    \item Context-aware vulnerability reasoning
    \item Automated remediation suggestions
    \item Cross-component security impact analysis
\end{itemize}

\subsection{Fine-Tuning Dataset Construction}

Drawing from Papers 5 and 9, we propose a systematic approach for constructing security-focused fine-tuning datasets:

\begin{enumerate}
    \item \textbf{Source Collection:} Aggregate vulnerability data from multiple sources including audit reports, CVE databases, and real-world incident reports
    \item \textbf{LLM-Assisted Extraction:} Use LLMs to extract structured vulnerability information from unstructured sources
    \item \textbf{CWE Classification:} Map vulnerabilities to the CWE hierarchy using tree-of-thought reasoning
    \item \textbf{Human Validation:} Implement systematic human validation with inter-rater agreement metrics
    \item \textbf{Continuous Evolution:} Design mechanisms for continuous dataset updates as new vulnerabilities emerge
\end{enumerate}

\subsection{Context-Aware Security Configuration}

Based on findings from Paper 10, we propose enhanced agent context files that explicitly incorporate security requirements:

\begin{itemize}
    \item \textbf{Security Sections:} Mandatory inclusion of security guidelines, acceptable patterns, and prohibited practices
    \item \textbf{Vulnerability Checklists:} Project-specific vulnerability prevention checklists
    \item \textbf{Security Testing Commands:} Explicit instructions for running security scans
    \item \textbf{Remediation Guidance:} Pre-approved remediation strategies for common issues
\end{itemize}

\subsection{Evaluation Metrics and Benchmarks}

We propose a standardized evaluation framework:

\begin{table}[H]
\centering
\caption{Proposed Evaluation Metrics}
\begin{tabular}{ll}
\toprule
\textbf{Metric Category} & \textbf{Specific Metrics} \\
\midrule
Detection Accuracy & Precision, Recall, F1-Score, MCC \\
Localization Accuracy & Line-level, Function-level, File-level precision \\
False Positive Rate & FP ratio, Alert fatigue index \\
Performance & Analysis latency, Throughput, Resource utilization \\
Coverage & CWE/OWASP category coverage, Tool coverage \\
\bottomrule
\end{tabular}
\end{table}

%==============================================================================
\section{Research Gaps and Future Directions}
%==============================================================================

\subsection{Identified Research Gaps}

Based on our analysis, we identify the following research gaps:

\begin{enumerate}
    \item \textbf{Security in Agent Context Files:} Paper 10 reveals that security requirements are specified in only 14.5\% of agent context files, indicating a critical gap in agentic coding practices.
    
    \item \textbf{Cross-Platform Security Consistency:} While GLITCH addresses polyglot detection, there remains a need for consistent security analysis across the entire DevOps toolchain.
    
    \item \textbf{Real-Time Security Feedback:} Current approaches focus on batch analysis; real-time security feedback during code development remains underexplored.
    
    \item \textbf{Automated Remediation Quality:} While LLMs can suggest remediations, the quality and security of automated fixes require further validation.
    
    \item \textbf{Domain-Specific Vulnerability Patterns:} Cloud-native, containerized, and serverless environments may have unique vulnerability patterns not fully captured by existing taxonomies.
\end{enumerate}

\subsection{Proposed Future Research Directions}

\subsubsection{Security-Aware LLM Fine-Tuning}
Develop comprehensive fine-tuning datasets and methodologies specifically designed for IaC and smart contract security, extending the work of GenSIaC and Forge.

\subsubsection{Hybrid Detection Frameworks}
Research optimal combinations of rule-based, ML-based, and LLM-based approaches for different use cases and performance requirements.

\subsubsection{Continuous Security Learning}
Develop mechanisms for security tools to continuously learn from new vulnerabilities and adapt detection patterns without manual intervention.

\subsubsection{Developer-Centric Security Tools}
Focus on improving developer experience with security tools, reducing false positive fatigue, and providing actionable remediation guidance.

%==============================================================================
\section{Conclusion}
%==============================================================================

This bibliographic report has analyzed ten significant research papers addressing security challenges in Infrastructure-as-Code and smart contracts. The analysis reveals a clear evolution from rule-based static analysis toward sophisticated LLM-based approaches, while highlighting the continued importance of traditional tools for high-assurance verification.

Key findings include:

\begin{enumerate}
    \item LLMs demonstrate superior recall compared to traditional static analysis tools but require careful management of false positives
    \item Semantic-aware approaches significantly improve detection accuracy
    \item Fine-tuning with domain-specific datasets substantially improves LLM performance for security tasks
    \item Security requirements are underrepresented in developer practices, particularly in agentic coding contexts
    \item A multi-tier approach combining the strengths of different methods appears most promising
\end{enumerate}

The proposed framework integrates rule-based detection, semantic-aware ML analysis, and LLM-assisted deep analysis into a comprehensive security pipeline. Future research should focus on security-aware LLM fine-tuning, hybrid detection frameworks, and developer-centric security tools.

%==============================================================================
\section*{Acknowledgments}
%==============================================================================

This bibliographic analysis was conducted to provide a comprehensive overview of the current state-of-the-art in security analysis for Infrastructure-as-Code and smart contracts.

%==============================================================================
\begin{thebibliography}{99}
%==============================================================================

\bibitem{war2025}
War, A., Diallo, A., Habib, A., Klein, J., Bissyandé, T.F. (2025). Vulnerabilities in Infrastructure as Code: What, How Many, and Who? \textit{Empirical Software Engineering}, 30:120.

\bibitem{saavedra2022}
Saavedra, N., Ferreira, J.F. (2022). GLITCH: Automated Polyglot Security Smell Detection in Infrastructure as Code. \textit{IEEE/ACM International Conference on Automated Software Engineering (ASE)}.

\bibitem{zheng2024}
Zheng, Z., Su, J., Chen, J., Lo, D., Zhong, Z., Ye, M. (2024). DAppSCAN: Building Large-Scale Datasets for Smart Contract Weaknesses in DApp Projects. \textit{IEEE Transactions on Software Engineering}, 50(6).

\bibitem{oliveira2025}
Oliveira, W., Paiva, F., Pereira, T.E., Brunet, J. (2025). A Defect Taxonomy for Infrastructure as Code: A Replication Study. \textit{arXiv:2505.01568v3}.

\bibitem{chen2026}
Chen, J., Shen, Y., Zhang, J., et al. (2026). Forge: An LLM-driven Framework for Large-Scale Smart Contract Vulnerability Dataset Construction. \textit{IEEE/ACM International Conference on Software Engineering (ICSE)}.

\bibitem{gnieciak2025}
Gnieciak, D., Szandala, T. (2025). Large Language Models Versus Static Code Analysis Tools: A Systematic Benchmark for Vulnerability Detection. \textit{IEEE Access}.

\bibitem{war2025b}
War, A., Nikiema, S.L.B., Samhi, J., Klein, J., Bissyandé, T.F. (2025). Security Smells in Infrastructure as Code Scripts: A Taxonomy Update Beyond the Seven Sins. \textit{arXiv:2509.18761v1}.

\bibitem{war2025c}
War, A., Rawass, A.A., Kabore, A.K., Samhi, J., Klein, J., Bissyandé, T.F. (2025). Detection of Security Smells in IaC Scripts Through Semantics-Aware Code and Language Processing. \textit{arXiv:2509.18790v1}.

\bibitem{li2025}
Li, Y., Grella, M., Nahmias, D., et al. (2025). GenSIaC: Toward Security-Aware Infrastructure-as-Code Generation with Large Language Models. \textit{arXiv:2511.12385v1}.

\bibitem{chatlatanagulchai2025}
Chatlatanagulchai, W., Li, H., Kashiwa, Y., et al. (2025). Agent READMEs: An Empirical Study of Context Files for Agentic Coding. \textit{arXiv:2511.12884v1}.

\end{thebibliography}

\end{document}
